% GW_Phase4_D1M2_Presentation
% Built from Docs/Formulation/self-consistent/SC_manuscript.
% Style follows GW_Phase4_D1M1_Presentation.
\documentclass[aspectratio=169,10pt]{beamer}

\usepackage[T1]{fontenc}
\usepackage{lmodern}
\usepackage{amsmath,amssymb,bm}
\usepackage{booktabs}
\usepackage{graphicx}
\usepackage{multirow}
\usepackage{array}

\usetheme{default}
\usefonttheme{professionalfonts}
\setbeamertemplate{navigation symbols}{}
\setbeamertemplate{blocks}[rounded][shadow=true]
\setbeamertemplate{itemize items}[circle]

\definecolor{gwblue}{RGB}{56,56,158}
\definecolor{gwlight}{RGB}{226,228,244}
\setbeamercolor{structure}{fg=gwblue}
\setbeamercolor{frametitle}{bg=gwblue,fg=white}
\setbeamercolor{title}{bg=gwblue,fg=white}
\setbeamercolor{block title}{bg=gwblue!75,fg=white}
\setbeamercolor{block body}{bg=gwlight!55,fg=black}
\setbeamerfont{frametitle}{size=\large,series=\mdseries}
\setbeamerfont{title}{size=\Large,series=\mdseries}

\setbeamertemplate{frametitle}{%
  \nointerlineskip%
  \begin{beamercolorbox}[wd=\paperwidth,ht=3.0ex,dp=1.4ex,leftskip=0.35cm]{frametitle}%
    \usebeamerfont{frametitle}\insertframetitle%
  \end{beamercolorbox}%
}
\setbeamertemplate{footline}{%
  \hfill\usebeamercolor[fg]{structure}\tiny
  \insertframenumber/\inserttotalframenumber\hspace{3mm}\vspace{2mm}%
}

\newcommand{\aloop}{$\langle a\rangle$}
\newcommand{\cloop}{$\langle c\rangle$}
\newcommand{\spvec}[1]{\bm{\mathsf{#1}}}
\newcommand{\spmat}[1]{\bm{\mathsf{#1}}}
\newcommand{\xv}{\bm{x}}
\newcommand{\FIG}{../Formulation/self-consistent/SC_manuscript/figures}
\newcommand{\foot}[1]{\par\vspace{1ex}%
  \begingroup\footnotesize\raggedright #1\par\endgroup}

\title{A Self-Consistent Spatially Resolved Cluster Dynamics Model\\for Irradiated Zirconium}
\author{N. Ghoniem, G. Po, M. Marron \& R. Escobar}
\date{August 31, 2026}

\begin{document}

% ---------------------------------------------------------------- title
{\setbeamertemplate{footline}{}
\begin{frame}[plain]
  \vfill
  \begin{beamercolorbox}[wd=\paperwidth,center]{title}
    \vspace{2.2ex}
    {\usebeamerfont{title}A Self-Consistent Spatially Resolved\\[4pt]
     Cluster Dynamics Model for Irradiated Zirconium}
    \vspace{2.2ex}
  \end{beamercolorbox}
  \vspace{5ex}
  \begin{center}
    {\large N. Ghoniem, G. Po, M. Marron \& R. Escobar}\\[3ex]
    {\footnotesize Multiscale Materials Solutions --- Deliverable D1/M2}\\[3ex]
    {\normalsize August 31, 2026}
  \end{center}
  \vfill
\end{frame}}

% ---------------------------------------------------------------- 2
\begin{frame}{What D1/M2 delivers}
\begin{columns}[T]
\begin{column}{0.48\textwidth}
\begin{block}{Three things}
\begin{itemize}\small
  \item A \textbf{self-consistent} cluster-dynamics model: every bias factor is
        generated from the \emph{same} diffusion tensors that transport the
        mobile species.
  \item A \textbf{continuum $\to$ discrete} conversion that preserves stored
        defects exactly and loop number to the integer draw.
  \item A \textbf{two-time-scale integration} that makes $2.1\times10^{8}$
        stiff equations tractable.
\end{itemize}
\end{block}
\end{column}
\begin{column}{0.48\textwidth}
\begin{block}{Demonstrated on}
\begin{itemize}\small
  \item A 500\,nm hexagonal $\alpha$-Zr single crystal, 800\,nm tall.
  \item $T=573$\,K, $G=10^{-7}\,\mathrm{dpa\,s^{-1}}$, zero applied stress.
  \item $0\to40$\,dpa in 13 dose intervals, 90\,663 nodes.
  \item Calibrated against \textbf{78 TEM measurements} across twelve
        temperatures and four decades of dose.
\end{itemize}
\end{block}
\end{column}
\end{columns}
\foot{The model is self-consistent in two senses --- in the bias factors, and
in the continuum/discrete transfer --- and both are load-bearing for what
follows.}
\end{frame}

% ---------------------------------------------------------------- 3
\begin{frame}{From D1/M1 to D1/M2}
\centering
\small\begin{tabular}{@{}p{0.24\textwidth}p{0.33\textwidth}p{0.35\textwidth}@{}}
\toprule
 & \textbf{D1/M1 (SRCD)} & \textbf{D1/M2 (this work)} \\
\midrule
Immobile families & 4 (aligned / non-aligned) & \textbf{9} crystallographic \\
Moments carried & density, content & \textbf{+ second moment $q$} \\
Capture efficiencies & phenomenological $Z=1+\delta$ & \textbf{Woo}, from the
  diffusion tensor \\
Size distribution & not resolved & \textbf{log-normal closure} \\
Grain boundary & Dirichlet, mobile only & \textbf{+ size-gated loop removal} \\
Discrete loops & visualization only & \textbf{conserving conversion} \\
Calibration & --- & \textbf{MAP against 78 measurements} \\
\bottomrule
\end{tabular}
\foot{The single structural change is that the anisotropy now enters
\emph{once}, through the diffusion tensors, and reaches the loops through the
capture efficiencies rather than through an independent set of fitted biases.}
\end{frame}

% ---------------------------------------------------------------- 4
\begin{frame}{The nine immobile families}
\centering
\begin{tabular}{@{}llll@{}}
\toprule
Family & Habit & Character & Burgers vector \\
\midrule
$a_1,a_2,a_3$ & prismatic $\{10\bar10\}$ & interstitial & $\tfrac13\langle11\bar20\rangle$ \\
$a_1^{v},a_2^{v},a_3^{v}$ & prismatic $\{10\bar10\}$ & vacancy & $\tfrac13\langle11\bar20\rangle$ \\
$c_0$ & basal, stacking-fault pyramid & vacancy & --- (no line) \\
$c_f$ & basal $(0001)$, faulted & vacancy & $\tfrac12[0001]$ \\
$c_p$ & basal $(0001)$, perfect & vacancy & $\tfrac12[0001]$ \\
\bottomrule
\end{tabular}
\vspace{1ex}
\foot{$|b_{\langle a\rangle}|=a=3.23$\,\AA{} and $|b_{\langle c\rangle}|=c/2=2.575$\,\AA{}
with the \emph{physical} axial ratio $c/a=1.5944$, not the ideal $\sqrt{8/3}$.
The same four numbers are used by the 0-D model, the continuum solver and the
discrete conversion --- a mismatch there was a factor of two in the \cloop{}
Burgers vector.}
\end{frame}

% ---------------------------------------------------------------- 5
\begin{frame}{Two populations, two kinds of equation}
\begin{columns}[T]
\begin{column}{0.48\textwidth}
\begin{block}{Mobile --- $v,i,2i,3i$}
\begin{itemize}\small
  \item Four concentration fields $\spvec c_M(\xv)$.
  \item \textbf{Spatially coupled} by diffusion.
  \item \textbf{No time derivative}: a steady reaction--diffusion problem.
  \item Dirichlet at grain boundaries.
  \item $4\times90\,663=362\,652$ unknowns.
\end{itemize}
\end{block}
\end{column}
\begin{column}{0.48\textwidth}
\begin{block}{Immobile --- nine loop families}
\begin{itemize}\small
  \item Three moments each: $n^k_{sL}$, $c^k_{sL}$, $q^k_{sL}$.
  \item \textbf{Pointwise}: no neighbor's state appears.
  \item Stiff ODEs in dose.
  \item $27\times90\,663=2\,447\,901$ unknowns.
\end{itemize}
\end{block}
\end{column}
\end{columns}
\foot{That asymmetry --- one block coupled and algebraic, the other pointwise
and stiff --- is what the operator split of \S4 exploits, and it is the reason
the calculation is affordable at all.}
\end{frame}

% ---------------------------------------------------------------- 6
\begin{frame}{The mobile master equation}
\begin{block}{}
\[
  \spvec 0 = \nabla\!\cdot\!\bigl(\spmat D_{\mathrm{ten}}\nabla\spvec c_M\bigr)
    + \spvec P_M
    + \tfrac12(\spmat I\otimes\spvec c_M^{T})\spmat R_2\spvec c_M
    + (\spmat E-\spmat D)\spvec c_M + \spvec S_{vL}
\]
\end{block}
\vspace{-1ex}
\begin{center}\small
\begin{tabular}{@{}ll@{}}
\toprule
$\nabla\!\cdot\!(\spmat D_{\mathrm{ten}}\nabla\spvec c_M)$ & anisotropic transport of the four mobile species \\
$\spvec P_M$ & cascade production (dose rate $\times$ survival $\times$ partition) \\
$\tfrac12(\spmat I\otimes\spvec c_M^{T})\spmat R_2\spvec c_M$ & second-order reactions: recombination and clustering \\
$\spmat E\spvec c_M$ & thermal dissociation of mobile clusters \\
$-\spmat D\spvec c_M$ & absorption at immobile sinks: network $+$ loops \\
$\spvec S_{vL}$ & vacancies returned by peripheral emission from loops \\
\bottomrule
\end{tabular}
\end{center}
\foot{Steady, not quasi-steady: $\partial_t\spvec c_M$ is absent by
construction, which is what allows the coupled problem to be solved a handful
of times over a march rather than at every step.}
\end{frame}

% ---------------------------------------------------------------- 7
\begin{frame}{Absorption: where the loops act back on the mobile species}
\[
  \spmat D = \mathrm{diag}\{\spvec\omega\odot\spvec C_s\}
             + \spmat D_d\,\spmat K_{\mathcal L},
  \qquad
  C_s^{v}=\frac{a^2}{z_c}\rho_N,\quad C_s^{i}=\frac{a^2}{z_c}Z_N\rho_N
\]
With $\omega^{m}=z_c\nu^{m}e^{-E^{m}_{\mathrm{mig}}/k_BT}$ and
$\overline D^{m}=(a^{2}/z_c)\,\omega^{m}$, the first term is simply
\[
  \spvec\omega\odot\spvec C_s
  =\bigl(\overline D^{v}\rho_N,\;Z_N\overline D^{i}\rho_N,\;
         Z_N\overline D^{2i}\rho_N,\;Z_N\overline D^{3i}\rho_N\bigr)
\]
\vspace{-1ex}
\begin{block}{}\small
$a^2$ and $z_c$ cancel: the network term is the classical
$\overline D^{m}k^{2}_N$. Splitting it into a frequency and a sink
concentration is 0-D bookkeeping, not a second physical quantity.
\end{block}
\foot{The loops enter the \emph{second} term only, through
$k_m^{2}=\sum_{(s,k)}Z_{sk,m}\rho^{k}_{sL}+\dots$ --- they must not also appear
in $\spvec C_s$, or their absorption would be counted once as a network sink
and once as loop growth.}
\end{frame}

% ---------------------------------------------------------------- 8
\begin{frame}{Anisotropy enters once, through one number per species}
\begin{columns}[T]
\begin{column}{0.50\textwidth}
\small
Migration energies are split at fixed $\overline D$:
\[
  E_m^{\langle11,22\rangle}=E_m^{\mathrm{eff}}+2k_BT\ln p^{m}
\]
\vspace{-2ex}
\[
  E_m^{\langle33\rangle}=E_m^{\mathrm{eff}}-4k_BT\ln p^{m}
\]
so a change in $p^{m}$ alters \emph{directionality}, never overall mobility.

\vspace{1ex}
The same $p^{m}$ builds the Woo capture efficiencies:
\[
  Z_{\mathrm{basal}}=Z^{0}_mp^{m},\qquad
  Z_{\mathrm{prism}}=Z^{0}_m\frac{p^{m}+p_m^{-2}}{2}
\]
\end{column}
\begin{column}{0.46\textwidth}
\begin{block}{Fixed from transport, not from loop counts}\small
\[
  p^{v}=1.02,\qquad p^{i}=0.70
\]
implying an interstitial migration anisotropy of \textbf{0.106\,eV} at 573\,K
--- a falsifiable prediction, open to atomistic test.
\end{block}
\end{column}
\end{columns}
\foot{This is the point of the self-consistency: because every bias follows
from the diffusion tensors, the anisotropy \emph{cannot} absorb a shortfall by
refitting a bias. It is an input, and what is left over has to appear somewhere
that can be named.}
\end{frame}

% ---------------------------------------------------------------- 9
\begin{frame}{Simultaneous \aloop{} and \cloop{} growth}
\small
Both growth conditions reduce to bounds on a single number --- the arrival ratio
\[
  A/B=\overline D^{v}c_v\Big/\textstyle\sum_m \overline D^{m}c_m|m|
\]
and the window between them has width $g(p^{v})/g(p^{i})$ with
$g(p)=2/(1+p^{-3})$, strictly increasing. Hence
\begin{block}{}
\centering\normalsize
Co-growth of both families is possible \textbf{if and only if} $p^{i}<p^{v}$.
\end{block}
\vspace{1ex}
The calibrated pair $p^{i}=0.70<p^{v}=1.02$ satisfies it with margin.
\foot{Two limits worth stating: the criterion governs the \emph{absorption flux}
only --- cascade nucleation is unaffected by anisotropy --- and an isotropic
calculation is another parameter point, not a zero of the criterion, so a ratio
taken against one cannot test a statement about signs.}
\end{frame}

% ---------------------------------------------------------------- 10
\begin{frame}{The immobile block: three moments per family}
\small
Restricting the master equation to indices that do not diffuse leaves a
birth--death equation on the integer size $m$. Three successive reductions:
\begin{enumerate}\small
  \item Treat $m$ as continuous, truncate Kramers--Moyal at second order
        $\Rightarrow$ Fokker--Planck in $m$.
  \item Multiply by $m^{j}$, integrate $\Rightarrow$ an \emph{exact} moment
        hierarchy; retain $j=0,1,2$.
  \item Close it with a log-normal $f_k(m)$.
\end{enumerate}
\[
  n^k_{sL}=\!\int\! f_k\,dm,\qquad
  c^k_{sL}=\!\int\! m f_k\,dm,\qquad
  q^k_{sL}=\!\int\! m^{2} f_k\,dm
\]
\[
  \bar m^k_{sL}=\frac{c^k_{sL}}{n^k_{sL}},\qquad
  \Delta_k\equiv\frac{q^k_{sL}n^k_{sL}}{(c^k_{sL})^{2}}=e^{s_k^{2}}\ \ge 1
\]
\foot{$\Delta_k\ge1$ is Cauchy--Schwarz on a non-negative measure. It is not a
tolerance: $\Delta<1$ is a defect in the equations, and it is the sharpest
available check that every channel is closed consistently.}
\end{frame}

% ---------------------------------------------------------------- 11
\begin{frame}{The second moment earns its place}
\small
The Fokker--Planck contribution to $\dot q$ is
\[
  \dot q^k_{sL}\big|_{\mathrm{FP}}
  = 2\Pi^k_{sL}\Bigl[\underbrace{\Sigma^{\mathrm{net}}_k\bar m^k_{sL}\Delta_k^{3/8}}_{\text{drift}}
    +\underbrace{\tfrac12\Sigma^{\mathrm{gro}}_k\Delta_k^{-1/8}}_{\text{spreading}}\Bigr]
\]
\begin{block}{}
At $\Sigma^{\mathrm{net}}_k=0$ the drift vanishes but
$\dot q^k_{sL}=\Pi^k_{sL}\Sigma^{\mathrm{gro}}_k\Delta_k^{-1/8}>0$: the
distribution \textbf{keeps broadening at fixed mean size and fixed number}.
No pair of lower moments can express that.
\end{block}
\vspace{0.5ex}
The closure reaches the rest of the model in exactly one place --- as
$\Delta_k^{-1/8}$ on the sink prefactor --- and reaches the discrete conversion
as the distribution loop sizes are drawn from.
\foot{Drift carries $\Delta^{3/8}$ and spreading $\Delta^{-1/8}$: a family that
has already broadened broadens \emph{faster}. The dispersion is self-amplifying
through the drift, which is a real feature and, at one node in the calculation,
a visible one.}
\end{frame}

% ---------------------------------------------------------------- 12
\begin{frame}{Loss channels, and the sign that is easy to get wrong}
\small
Coalescence, by an Avrami overlap on the family's own size:
\[
  \varphi_{LL}=1-e^{-\kappa_{LL}\frac43\pi R^3_{LL}N},\qquad
  \varphi_{LN}=1-e^{-\kappa_{LN}\pi R^2_{LN}\rho_N}
\]
\[
  \Delta N=\nu_{LL}\varphi_{LL}n+\Delta N_{LN},\qquad
  \Delta Q=\Delta_k\bigl(\bar m^{2}\Delta N-2\bar m\Delta C\bigr)
\]
\begin{block}{A coalescing loop does not leave the family --- it MERGES}
Its defects stay; only $\Delta C_{LN}$ is lost. Fewer loops holding the same
content are \emph{larger} loops, so like-loop coalescence \textbf{raises} $q$.
Debiting $\Delta N\langle m^2\rangle$ --- the natural-looking reading ---
inverts the sign on the largest single term in the equation.
\end{block}
\foot{Measured consequence of getting it wrong: $q$ crossed zero at 1\,dpa and
was negative on 40\% of the $c_p$ nodes at 10\,dpa. Coalescence carries
$\sim98\%$ of the loop-number loss, so this is not a small term.}
\end{frame}

% ---------------------------------------------------------------- 13
\begin{frame}{The distribution leaks at both ends}
\small
\begin{center}
\begin{tabular}{@{}lll@{}}
\toprule
 & \textbf{Small end} (vanishing) & \textbf{Large end} (coalescence) \\
\midrule
Removes & loops below $m_{\mathrm{van}}$ & the biggest loops \\
Weight & $1-\Phi^{(j)}$, \emph{below} the mean & $\Phi^{(j)}$, \emph{above} it \\
Effect on $\bar m$ & \textbf{raises} it & lowers it \\
Realizability & automatic --- what remains is a genuine sub-measure \\
\bottomrule
\end{tabular}
\end{center}
\[
  \Xi^k_j=\bigl[1-\Phi^{(j)}_k(m_{\mathrm{van}})\bigr]\times
  \{n^k_{sL},\,c^k_{sL},\,q^k_{sL}\},\qquad
  \Phi^{(j)}_k(m_\theta)=\tfrac12\mathrm{erfc}\!\left(
    \frac{\ln m_\theta-\mu_k-js_k^{2}}{s_k\sqrt2}\right)
\]
\begin{block}{}
Because $\Phi^{(2)}>\Phi^{(1)}>\Phi^{(0)}$, the complements reverse that
ordering: what the small-end leak removes is \emph{smaller} than the family
mean, so \textbf{the density falls while the surviving mean size rises}.
\end{block}
\foot{That is the one trade coalescence cannot make. It was the first channel
found that moves density and diameter the way the measurements want them
moved --- fewer \emph{and} larger.}
\end{frame}

% ---------------------------------------------------------------- 14
\begin{frame}{Grain boundaries: two populations, two mechanisms}
\begin{columns}[T]
\begin{column}{0.46\textwidth}
\begin{block}{Mobile --- Dirichlet}\small
\[c^{v}=e^{-E^{f}_v/k_BT},\quad c^{i}=c^{2i}=c^{3i}=0\]
on $\partial\Omega_D$. Thermal equilibrium is imposed; the layer that develops
is a result, not an input.
\end{block}
\end{column}
\begin{column}{0.50\textwidth}
\begin{block}{Loops --- size-gated removal}\small
A loop of radius $R^k=\lambda_k\sqrt m$ at distance $x$ touches the face when
\[m\ge m_{gb}(x)=(x/\lambda_k)^{2}\]
so the absorbed fraction of each moment is
$\Theta^k_j=\Phi^{(j)}_k(m_{gb})\times\{n,c,q\}$.
\end{block}
\end{column}
\end{columns}
\vspace{1ex}
\foot{A loop cannot be given a boundary condition --- it is untransported and
has no flux to set. What it has is a \emph{size}, and the family already
carries the distribution of sizes needed to say what fraction of it reaches the
wall. The channel costs no new parameter beyond a rate $\nu_{gb}$.}
\end{frame}

% ---------------------------------------------------------------- 15
\begin{frame}{The operator split, and why it is legitimate}
\small
Lie--Trotter, alternating over each substep:
\begin{enumerate}\small
  \item \textbf{Fast solve} --- the steady mobile field
        $\spvec c_M^{*}(\xv)$ for the immobile state currently held.
  \item \textbf{Slow march} --- the immobile ODEs integrated at every point
        with the mobile species \emph{frozen}.
\end{enumerate}
\begin{block}{The decoupling is by construction, not by approximation}
The only thing coupling neighboring points is diffusion of the mobile species.
Freeze the mobile field and the spatial coupling is \emph{gone}: a point's
immobile ODEs depend on its own loop populations and the four numbers pinned
into $\spvec c_M$.
\end{block}
\vspace{0.5ex}
Splitting error is first order in the dose spanned between fast solves --- so
the fast-solve cadence, not the substep count, is the accuracy dial.
\foot{The bill is paid in splitting error and nowhere else. The coupled problem
has not been approximated away; it has been moved to where it is affordable.}
\end{frame}

% ---------------------------------------------------------------- 16
\begin{frame}{What the split buys, in numbers}
\centering
\begin{tabular}{@{}lr@{}}
\toprule
 & flops for one implicit factorization \\
\midrule
Split: $73\,518$ independent $36\times36$ systems & $3.4\times10^{9}$ \\
Coupled: one $3\,263\,868$-dimensional system & $3.5\times10^{19}$ \\
\midrule
\textbf{Ratio} & $\bm{1.0\times10^{10}}$ \\
\bottomrule
\end{tabular}
\vspace{2ex}

\begin{columns}[T]
\begin{column}{0.47\textwidth}
\begin{block}{Measured, 500\,nm case}\small
\begin{tabular}{@{}lr@{}}
march total & 1\,h 10\,m \\
one fast solve ($\times26$) & 1\,m 54\,s \\
one slow substep ($\times78$) & 15.3\,s \\
dedup factor & $\times1.23$ \\
\end{tabular}
\end{block}
\end{column}
\begin{column}{0.47\textwidth}
\begin{block}{Why frequency matters}\small
The fast solve has \emph{no time derivative}, so the coupled problem is solved
\textbf{26 times over the march} rather than at every timestep. The
counterfactual is $\sim$25\,h for the 4-species mobile block alone.
\end{block}
\end{column}
\end{columns}
\foot{Sparsity softens the coupled figure but does not rescue it: a 3-D FE
Jacobian still factorizes at $O(N^2)$ with fill growing as $O(N^{4/3})$.}
\end{frame}

% ---------------------------------------------------------------- 17
\begin{frame}{Continuum $\to$ discrete: what must be preserved}
\small
\begin{center}
\begin{tabular}{@{}lll@{}}
\toprule
Invariant across a transfer & Requirement & Achieved \\
\midrule
Stored defects & exact & to $10^{-16}$ (rescaled) \\
Loop number & exact & to the integer draw \\
Size dispersion & preserved & to the sampling truncation \\
Sink strength & continuous & 0.97--1.04 at every dose \\
\bottomrule
\end{tabular}
\end{center}
\[
  \lambda^k_{sj}=\frac{n^k_{sL}(\xv_j)}{\Omega}V_j,\qquad
  \Lambda^k_s=\sum_j\lambda^k_{sj},\qquad
  N_{\mathrm{draw}}=\mathrm{round}(\Lambda^k_s)
\]
\begin{block}{}
A calculation may therefore be run as continuum, as discrete, or as a hybrid,
with the choice made on grounds of \textbf{dose and cost} rather than of
physics.
\end{block}
\foot{Rounding rather than drawing a Poisson total is deliberate: $\mathrm{Var}
=\Lambda$ there, and a dose sweep whose count fluctuates by $\sqrt\Lambda$
between snapshots is dominated by counting noise.}
\end{frame}

% ---------------------------------------------------------------- 18
\begin{frame}{Nodal volumes, and what a sparse cell means}
\small
$V_j$ is estimated by \textbf{Monte Carlo}: points drawn uniformly in the
bounding box, rejected against the convex hull, each survivor assigned to its
nearest node.
\[
  V_j=\frac{M_j}{M}V_\Omega,\qquad M=4\times10^{6},\qquad
  \text{error}\sim M^{-1/2}
\]
\begin{columns}[T]
\begin{column}{0.48\textwidth}
\begin{block}{The one condition}\small
The cells must tile the crystal exactly. They do: every accepted sample is
charged to exactly one node, so $\sum_jV_j=V_\Omega$ identically, however
coarse the sampling. Error moves volume \emph{between} cells; it cannot change
their sum.
\end{block}
\end{column}
\begin{column}{0.48\textwidth}
\begin{block}{Two things it is not}\small
Not the FEM mesh --- the cells are built on the concentration trial-function
nodes and share only the point set. And a node with \emph{zero} volume is a
correct answer, not a failure: $\sim12\%$ receive none.
\end{block}
\end{column}
\end{columns}
\vspace{0.5ex}
\foot{No cell expects a whole loop: $\lambda\sim10^{-2}$ typically. The draw is
taken once for the \emph{family} and the nodes sampled from the multinomial ---
rounding node by node would return zero everywhere and convert nothing.}
\end{frame}

% ---------------------------------------------------------------- 19
\begin{frame}{Calibration: what is fitted and what is not}
\small
\begin{center}
\begin{tabular}{@{}lll@{}}
\toprule
Class & Examples & Fitted? \\
\midrule
Model controls & which families, which channels & no --- selects physics \\
Crystallography & $a$, $c$, $\Omega$, Burgers vectors & no --- measured \\
Transport & $p^{v}=1.02$, $p^{i}=0.70$ & \textbf{no --- pinned} \\
Rate parameters & capture prefactors, coalescence & \textbf{yes} \\
\bottomrule
\end{tabular}
\end{center}
\vspace{1ex}
\textbf{Targets:} 78 TEM measurements --- 73 \aloop{} and 5 \cloop{}, twelve
temperatures, four decades of dose.

\textbf{Estimator:} maximum \emph{a posteriori} under stated priors, on
$\log_{10}(\text{model}/\text{experiment})$, symmetric under inversion of the
ratio.

\textbf{Search:} staged bounded Nelder--Mead, polished jointly. One objective
evaluation is 0.2--0.5\,s, which is what makes a refit tractable.
\foot{The 0-D estimator and the 3-D slow step are the \emph{same binary with
the same switches}, differing only in whether the mobile species are frozen ---
so a parameter set fitted in 0-D is one the 3-D code runs term for term.}
\end{frame}

% ---------------------------------------------------------------- 20
\begin{frame}{Prismatic \aloop{} loops against experiment}
\begin{columns}[c]
\begin{column}{0.5\textwidth}
\centering\includegraphics[width=\textwidth]{\FIG/fit_A_density.png}
\end{column}
\begin{column}{0.5\textwidth}
\centering\includegraphics[width=\textwidth]{\FIG/fit_A_diameter.png}
\end{column}
\end{columns}
\foot{Number density (left) and diameter (right) by temperature. Density
$0.88\times$ the measurement over 71 points and diameter $1.21\times$ over 44,
across twelve temperatures and four decades of dose.}
\end{frame}

% ---------------------------------------------------------------- 21
\begin{frame}{Basal \cloop{} loops, and the parity plot}
\begin{columns}[c]
\begin{column}{0.34\textwidth}
\centering\includegraphics[width=\textwidth]{\FIG/fit_C_density.png}
\end{column}
\begin{column}{0.34\textwidth}
\centering\includegraphics[width=\textwidth]{\FIG/fit_C_diameter.png}
\end{column}
\begin{column}{0.30\textwidth}
\centering\includegraphics[width=\textwidth]{\FIG/fit_parity.png}
\end{column}
\end{columns}
\foot{On the \emph{fixed} anisotropy the calibration reproduces the basal loop
diameter without it being adjusted --- \textbf{100.7\,nm against 95\,nm
measured}. What it does not reach is the basal density, which remains
$2.1\times$ the measurement.}
\end{frame}

% ---------------------------------------------------------------- 22
\begin{frame}{What the calibration does not achieve}
\begin{columns}[T]
\begin{column}{0.48\textwidth}
\begin{block}{A prismatic temperature trend}\small
The mean conceals it: resolved by temperature the model runs high at both ends
and low in the middle, a factor of six low near 600\,K. Four independent
searches reach the same place.
\vspace{0.5ex}

Two readings the data cannot separate: the nucleation term may carry the wrong
temperature dependence, or the high-$T$ points may not measure the same
population.
\end{block}
\end{column}
\begin{column}{0.48\textwidth}
\begin{block}{Four parameters on their bounds}\small
$Z^{0}_{v}=1.30$ and $\xi_c=2.00$ at their ceilings; both coalescence exponents
at theirs.
\vspace{0.5ex}

$Z^{0}_{v}$ and $\xi_c$ both multiply the basal capture rate and are nearly
degenerate --- one statement, not two: \textbf{given a stated anisotropy the
model cannot deliver the measured basal population without a prefactor no
transport argument supplies.}
\end{block}
\end{column}
\end{columns}
\foot{That is sharper than earlier calibrations supported. With $p^{v}$ free it
absorbed the discrepancy; pinned, the shortfall is forced into the capture
prefactor and \emph{not} the diffusion tensor. It is testable by
character-resolved measurements at 600--700\,K.}
\end{frame}

% ---------------------------------------------------------------- 23
\begin{frame}{The demonstration case}
\begin{columns}[c]
\begin{column}{0.46\textwidth}
\centering\includegraphics[width=0.92\textwidth]{\FIG/fe_mesh.png}
\end{column}
\begin{column}{0.50\textwidth}
\small
\begin{tabular}{@{}lr@{}}
\toprule
Geometry & hexagonal prism \\
Width $\times$ height & 500 $\times$ 800\,nm \\
Elements & 60\,965 tets \\
CD nodes & 90\,663 \\
Boundary layer & 100\,nm at 20\,nm \\
\midrule
Temperature & 573\,K \\
Dose rate & $10^{-7}\,\mathrm{dpa\,s^{-1}}$ \\
Applied stress & zero \\
Dose range & $10^{-4}\to40$\,dpa \\
Intervals & 13 \\
\bottomrule
\end{tabular}
\end{column}
\end{columns}
\foot{Dirichlet on every face, so the entire surface is grain boundary. The
mesh is cut in half so the interior refinement is visible.}
\end{frame}

% ---------------------------------------------------------------- 24
\begin{frame}{Mobile species, early and late}
\centering\includegraphics[height=0.66\textheight]{\FIG/mobile_early_late.png}
\foot{$10^{-4}$\,dpa (left) and 40\,dpa (right), on two orthogonal mid-cuts.
The boundary layers order themselves by species mobility --- the vacancy, least
mobile, is depleted over the shortest distance --- and narrow by a factor of two
as $k^{2}$ grows with the loop population the same calculation is building.}
\end{frame}

% ---------------------------------------------------------------- 25
\begin{frame}{The three moments: basal faulted family $c_f$}
\centering\includegraphics[height=0.66\textheight]{\FIG/moments_cf.png}
\foot{Number density, content and second moment at $10^{-4}$ and 40\,dpa. The
three rows are not three views of one thing --- they answer different questions
and their spatial structures differ, which is the argument for carrying all
three.}
\end{frame}

% ---------------------------------------------------------------- 26
\begin{frame}{The three moments: prismatic variant $a_1$}
\centering\includegraphics[height=0.66\textheight]{\FIG/moments_a1.png}
\foot{The same three moments for one prismatic interstitial variant. The
prismatic family stays small and numerous where the basal family grows large and
sparse --- the two are not scaled versions of one another.}
\end{frame}

% ---------------------------------------------------------------- 27
\begin{frame}{The atom balance closes --- on an open system}
\begin{columns}[c]
\begin{column}{0.5\textwidth}
\centering\includegraphics[width=\textwidth]{\FIG/conservation_v.png}
\end{column}
\begin{column}{0.5\textwidth}
\centering\includegraphics[width=\textwidth]{\FIG/fractions_v.png}
\end{column}
\end{columns}
\foot{The six accumulators integrate a 0-D balance with \emph{no transport
term}, so on a Dirichlet domain the residual is not solver error --- it is the
atoms that left through the surface. At 40\,dpa the boundary absorbs
\textbf{96\%} of interstitial production as mobile species and a further
\textbf{27\%} of vacancy production as loops.}
\end{frame}

% ---------------------------------------------------------------- 28
\begin{frame}{The loop-denuded zone}
\centering\includegraphics[width=0.94\textwidth]{\FIG/denuded_zone.png}
\foot{Solid: grain-boundary loop absorption active. Dashed: the identical
calculation without it, differing in that one switch and nothing else. The zone
is $\sim5$\,nm wide, and across it the mean loop diameter recovers from
\textbf{6.8 to 144\,nm}. The channel also \emph{raises} the \cloop{} density in
the deep interior by a factor of 4--5, because removing the shell's sinks
changes the mobile field everywhere --- something a mean-field $k^2_{gb}$
structurally cannot do.}
\end{frame}

% ---------------------------------------------------------------- 29
\begin{frame}{Domain mean or interior mean? It is not a detail}
\centering\small
\begin{tabular}{@{}llrrrr@{}}
\toprule
& Region & $N_a$ (m$^{-3}$) & $d_a$ (nm) & $N_c$ (m$^{-3}$) & $d_c$ (nm) \\
\midrule
\multirow{2}{*}{loop channel off}
 & domain   & $3.79\times10^{23}$ & 4.29 & $3.98\times10^{23}$ & 17.60 \\
 & interior & $2.80\times10^{21}$ & 5.79 & $1.75\times10^{21}$ & 249.84 \\
\midrule
\multirow{2}{*}{loop channel on}
 & domain   & $3.86\times10^{21}$ & 5.58 & $2.26\times10^{21}$ & 103.79 \\
 & interior & $7.06\times10^{21}$ & 5.98 & $5.31\times10^{21}$ & 133.01 \\
\bottomrule
\end{tabular}
\vspace{1ex}
\foot{At 40\,dpa. \textbf{Without} the channel the two differ by $135\times$ in
$N_a$ and $227\times$ in $N_c$ --- the domain average is measuring the boundary
shell and nothing else. \textbf{With} it the disparity collapses to $1.8\times$
and $2.4\times$ and changes sign, the interior now carrying more loops than the
domain mean. That is the strongest single argument for the channel: it removes
an artifact that makes the most natural summary statistic meaningless.}
\end{frame}

% ---------------------------------------------------------------- 30
\begin{frame}{Loop density and size against dose}
\begin{columns}[c]
\begin{column}{0.5\textwidth}
\centering\includegraphics[width=\textwidth]{\FIG/loop_density.png}
\end{column}
\begin{column}{0.5\textwidth}
\centering\includegraphics[width=\textwidth]{\FIG/loop_sizes_vs_exp.png}
\end{column}
\end{columns}
\foot{Model lines against the measured points, $10^{-4}\to40$\,dpa. The march
now reaches beyond the \cloop{} measurements at 35\,dpa, so the comparison is at
matched dose rather than extrapolated. Everything saturates by $\sim$1\,dpa:
from 1 to 40\,dpa the densities move under $9\%$ and the radii not at all.}
\end{frame}

% ---------------------------------------------------------------- 31
\begin{frame}{Against experiment --- and the direction of the discrepancy}
\small
At 573\,K and 24.5\,dpa the measurements give
$N_c=8.0\times10^{20}\,\mathrm{m^{-3}}$ and $d_c=95$\,nm. The interior at
40\,dpa:
\begin{center}
\begin{tabular}{@{}lrr@{}}
\toprule
 & $N_c/N^{\exp}$ & $d_c/d^{\exp}$ \\
\midrule
loop channel on & 6.6 & 1.40 \\
loop channel off & 2.2 & 2.63 \\
\bottomrule
\end{tabular}
\end{center}
\begin{block}{The direction has changed with the refit}
Earlier parameter sets produced \emph{more, smaller} \cloop{} loops --- the
signature of a model removing loops faster than it grows them. This one
produces \textbf{more and larger}: stored basal content, which scales as
$N_cd_c^{3}$, is $18\times$ what is observed.
\end{block}
\foot{The two errors no longer oppose each other, and that is a harder failure
to excuse. It is not a size problem a growth term could absorb --- it is too
much vacancy content retained in basal loops.}
\end{frame}

% ---------------------------------------------------------------- 32
\begin{frame}{The conversion, exercised at four doses}
\centering\small
\begin{tabular}{@{}lrrrr@{}}
\toprule
Family & $10^{-4}$\,dpa & $10^{-2}$\,dpa & 1\,dpa & 40\,dpa \\
\midrule
$c_f$ (basal, vacancy)          & 5 & 238 & 642 & 691 \\
$a_1$ (prismatic, interstitial) & 1 &  55 & 213 & 232 \\
$a_1^{v}$ (prismatic, vacancy)  & 1 &  23 &  69 &  75 \\
\bottomrule
\end{tabular}
\vspace{1.5ex}

The number of loops is not a parameter of the conversion: it is
$\sum_jn_jV_j$, the number the field says is present. The ledger closes on
stored defects to machine precision and on loop number to the integer draw at
every one of the four doses.
\foot{The three interstitial prismatic variants stay \emph{exactly} equal in
number (232 each at 40\,dpa), which is the isotropic variant weighting under
zero applied stress; the three vacancy variants likewise (75 each).}
\end{frame}

% ---------------------------------------------------------------- 33
\begin{frame}{The discrete basal population}
\centering\includegraphics[height=0.64\textheight]{\FIG/loops_c_4dose.png}
\foot{(a) $10^{-4}$, (b) $10^{-2}$, (c) 1 and (d) 40\,dpa, drawn at true radius
over the continuum content field. Loops are clipped at the crystal surface for
drawing; the exported loops are whole.}
\end{frame}

% ---------------------------------------------------------------- 34
\begin{frame}{Every family superimposed}
\centering\includegraphics[width=\textwidth]{\FIG/loops_all_4dose.png}
\foot{Basal \cloop{} in blue, the interstitial prismatic variants and the
vacancy variants in their own colours, at true radius and sampled position. What
the sequence shows is not the basal family becoming more numerous but its loops
becoming very much \emph{larger}, until at 40\,dpa they dominate the drawing at
a diameter an order of magnitude above the prismatic variants.}
\end{frame}

% ---------------------------------------------------------------- 35
\begin{frame}{Two regimes, and the limits of the mean field}
\begin{columns}[T]
\begin{column}{0.48\textwidth}
\begin{block}{At $10^{-4}$\,dpa}\small
5 basal and 6 prismatic loops in
$1.3\times10^{8}\,\mathrm{nm^{3}}$ --- a population so dilute that a continuum
density describing it is a statement about an \emph{ensemble}, not about this
crystal.
\end{block}
\end{column}
\begin{column}{0.48\textwidth}
\begin{block}{At 40\,dpa}\small
Basal loops 131\,nm across at a mean spacing of 57\,nm --- each loop is
\emph{wider than the distance to its neighbour}. This is where the Avrami
kernels work hardest and are least defensible.
\end{block}
\end{column}
\end{columns}
\vspace{1ex}
\begin{block}{And most basal loops do not fit}
At 40\,dpa, \textbf{419 of 691} \cloop{} loops (60.6\%) are centered within
their own radius of a face and overhang it, by up to 88\,nm; the prismatic
figure is 3.5\%. A discrete-dislocation solver would decline them, and their
defects would remain in the continuum.
\end{block}
\foot{The conversion spans both regimes, which is what makes it a check on the
mean field rather than a convenience --- and the refusal fraction is the
argument for a larger crystal.}
\end{frame}

% ---------------------------------------------------------------- 36
\begin{frame}{Reconstructed size distributions}
\centering\includegraphics[height=0.62\textheight]{\FIG/size_spectrum_dose.png}
\foot{Population per unit $\ln d$ per cubic metre at four doses, over six
decades below the peak. By 1\,dpa the \emph{shape} has essentially stopped
changing: the 1\,dpa and 40\,dpa curves are nearly coincident in every panel, so
the population saturates in shape as well as in its aggregates.}
\end{frame}

% ---------------------------------------------------------------- 37
\begin{frame}{One node misbehaves, and it is structural}
\small
At 1 and 40\,dpa a single node of the 90\,663 --- node 51\,609, at
$x=214.2$\,nm, the point of the prism \emph{furthest} from any face --- carries
$\Delta=4.7\times10^{3}$ and $6.2\times10^{3}$, against a median of 1.070.
Every other node stays below $\Delta=1.5$.
\begin{block}{The mechanism is in the moment equation}
Drift carries $\Delta_k^{3/8}$ and spreading $\Delta_k^{-1/8}$: a family that
has already broadened broadens \emph{faster} while its spreading is damped ---
the dispersion is self-amplifying through the drift. Everywhere else that
feedback is held in check by coalescence and by the finite mobile flux; at the
one point where the flux is largest and no boundary channel acts, it is not.
\end{block}
\foot{$\Delta\ge1$ is never violated, so no realizability check fires. The
remedy is not a tolerance --- it is a statement about where the closure is
valid, and the gap between the outlier and the rest of the crystal is three
orders of magnitude with nothing in between.}
\end{frame}

% ---------------------------------------------------------------- 38
\begin{frame}{Conclusions}
\begin{columns}[T]
\begin{column}{0.48\textwidth}
\begin{block}{What is established}\small
\begin{itemize}
  \item Every bias follows from the diffusion tensors through one $p^{m}$ per
        species. The anisotropy is an \textbf{input}, not a fitted absorber.
  \item Spatial resolution matters: boundary layers order by mobility and
        narrow by $2\times$; a domain mean without the loop channel measures
        the boundary shell.
  \item The transfer preserves what defines the population, so
        continuum / discrete / hybrid is a \textbf{cost} decision.
  \item $2.1\times10^{8}$ stiff equations in 1\,h 10\,m on 8 cores.
\end{itemize}
\end{block}
\end{column}
\begin{column}{0.48\textwidth}
\begin{block}{What is open}\small
\begin{itemize}
  \item The basal density, $2.1\times$ the measurement, with the capture
        prefactor on its bound --- absorption, not transport.
  \item The prismatic temperature trend, a factor of six low near 600\,K.
  \item $p^{i}=0.70$ implies $0.106$\,eV of migration anisotropy at 573\,K ---
        \textbf{open to atomistic test}.
  \item Whether $c_p$ needs a sink of its own; it is a terminal population.
\end{itemize}
\end{block}
\end{column}
\end{columns}
\foot{The useful outcome is that a shortfall now has to appear somewhere it can
be named. With the anisotropy pinned, the model cannot quietly refit a bias to
hide it.}
\end{frame}

\end{document}
